\chapter{Background}
\label{chap:background}

% Self-contained theory chapter ("self-study" material). Sections follow the
% light: each concept is introduced when the previous one creates the need
% for it. Populate each section below; reorder/add as you see fit.

\section{Radiometry}
\label{sec:bg_radiometry}
% On-ramp: flux, irradiance, radiance, solid angle. The quantities all later
% equations are written in.

\section{The Rendering Equation and Global Illumination}
\label{sec:bg_rendering_equation}
% Surface light transport: BRDF, the rendering equation, why solving it
% globally is hard. Reflection and refraction (Fresnel, Snell) as the BSDF
% cases that later give rise to caustics.

\section{Light Transport in Participating Media}
\label{sec:bg_media}
% Extend to volumes: absorption, out-/in-scattering, emission; extinction
% and albedo; phase functions (isotropic, Henyey-Greenstein); transmittance
% and the volume rendering equation (RTE).

\section{Monte Carlo Integration}
\label{sec:bg_monte_carlo}
% Estimators, variance, importance sampling; distance/free-flight sampling
% for media; path tracing as the canonical unbiased solver and why it is
% slow for media.

\section{Density Estimation and Photon Mapping}
\label{sec:bg_density_estimation}
% The bridge from MC to our method: two-pass light transport, photon maps,
% kernel density estimation, radiance estimate on surfaces and in media;
% bias/consistency trade-off. Historical variants stay in Related Work.

\section{Ray Tracing and Rasterization}
\label{sec:bg_raster_vs_rt}
% The two computational paradigms we hybridize: visibility by ray casting
% vs. by projection/scan conversion; the GPU pipeline, billboards/splatting,
% blending -- what rasterization gives us for free and what it cannot do.
